\section{TriGAT --- Triplet-Based Primitive}
\label{sec_meth_trigat}

This section describes a variant of PG-GAT in which the per-node primitive is replaced from a single joint to a skeleton triplet of three connected joints.
The motivation is articulation:
two people standing in similar absolute positions remain distinguishable by their joint angles
(shoulder--elbow--wrist, hip--knee--ankle, head--torso),
and the pivot-angle of a triplet is rotation-invariant articulation information that single-joint embeddings do not directly capture.
This intuition was converted into the experimental hypothesis
that embedding triplets rather than joints would surface articulation as a first-class signal
and thereby produce a more discriminative embedding for the hard case of spatially overlapping people.
The result reported in Chapter~\ref{chap_4} contradicts this hypothesis;
this section describes the methodology in enough detail that the negative result is interpretable.

\subsection{Triplet enumeration and features}
\label{sec_meth_trigat_features}

The COCO skeleton graph admits exactly $19$ canonical triplets,
indexed by joint-type triples $(t_A, t_B, t_C)$ where $(t_A, t_B)$ and $(t_B, t_C)$ are both COCO skeleton edges.
The pivot is the middle joint $B$.
For each scene, every detected person contributes at most $19$ triplets;
triplets with a missing or invisible joint are skipped,
which leaves the per-person triplet count between roughly $12$ and $19$ depending on visibility.

A triplet node carries a continuous feature vector
\begin{equation}
    \mathbf{f}_{ABC} = \bigl[\,\mathbf{p}_B,\;\Delta_{AB},\;\Delta_{CB},\;\cos\theta,\;\sin\theta,\;\|\Delta_{AB}\|,\;\|\Delta_{CB}\|\,\bigr] \in \mathbb{R}^{10},
    \label{eq:trigat_features}
\end{equation}
where $\mathbf{p}_B \in [0,1]^{2}$ is the pivot position,
$\Delta_{AB} = \mathbf{p}_A - \mathbf{p}_B$ and $\Delta_{CB} = \mathbf{p}_C - \mathbf{p}_B$ are the two outgoing bones,
and $\theta$ is the pivot angle $\angle(A, B, C)$ measured by the cosine and sine of its values
(which avoids the wrap-around discontinuity of representing the angle directly).
A learnable $19$-way triplet-type embedding $\mathbf{e}^{\text{tri}}_{(t_A, t_B, t_C)} \in \mathbb{R}^{16}$ is concatenated with $\mathbf{f}_{ABC}$ to form the layer-zero node feature.
The triplet-type embedding plays the same role for triplets that the joint-type embedding of Section~\ref{sec_meth_problem_graph} plays for joints.

\subsection{Triplet graph and embedding network}
\label{sec_meth_trigat_arch}

The triplet graph $\mathcal{G}^{\text{tri}} = (\mathcal{V}^{\text{tri}}, \mathcal{E}^{\text{tri}})$
has one node per detected triplet,
so $|\mathcal{V}^{\text{tri}}|$ is roughly $19K$ for a $K$-person scene with full visibility
(comparable in scale to the joint graph's $17K$).
Edges are formed by $k$-nearest neighbours in pivot-position space with $k = 8$,
matching the convention of the joint graph.
The embedding network on top of this graph is a GATv2 stack with the same depth, width, and head count as the PG-GAT default,
and the same contrastive loss of Equation~\eqref{eq:contrastive_loss} is applied at the triplet level
with the supervisory label set to the pivot joint's person identity.
The per-triplet output is a $128$-dimensional embedding on the unit hypersphere.

\subsection{Voting from triplet clusters to joint labels}
\label{sec_meth_trigat_voting}

At inference, $k$-means with $k = K$ is applied to the triplet embeddings,
which yields a per-triplet cluster label.
Each joint participates in up to three triplets,
each of which has been assigned an arbitrary cluster identifier by $k$-means,
and a deterministic voting rule converts triplet-cluster labels into joint labels.
The vote weights the pivot triplet at twice the weight of the wing triplets
on the basis that the pivot articulation is the most distinguishing feature of any triplet.
Ties are broken by the lowest-numbered triplet's cluster.
A Hungarian matching of cluster identifiers to ground-truth person identifiers follows
to compute the final PGA score per Equation~\eqref{eq:pga}.

The voting step is the methodological soft spot of the TriGAT design,
and the Chapter~\ref{chap_4} results identify it as the dominant failure mode:
a single mis-assigned triplet in a person can move two or three of that person's joints to a wrong cluster,
amplifying small triplet-level errors into joint-level mistakes.
The discussion in Section~\ref{sec_disc_trigat} treats the voting step in detail.

\subsection{Training schedule}
\label{sec_meth_trigat_training}

TriGAT is trained for 50 epochs on the synthetic training split with the optimiser settings of Section~\ref{sec_meth_baseline_training}.
The training budget matches PG-GAT exactly so that the comparison reported in Chapter~\ref{chap_4} controls for compute.
At inference, $k$-means with $k = K$ is applied to the triplet embeddings
and the resulting partition is converted to joint labels by the voting rule of Section~\ref{sec_meth_trigat_voting}.
Constrained-clustering read-outs of the same embeddings exist in companion work and are not reported here per the scoping rule of Section~\ref{sec_meth_problem_subproblems}.
